% Fuentes y límites de redacción: MAPA_FUENTES.md, capítulo 2.
\chapter{Contexto administrativo y fuente de información}
\label{ch:contexto}

\section{Tramitación de proyectos de generación renovable}
La tramitación de una instalación de generación reúne decisiones sobre su
conexión a la red, sus efectos ambientales y su autorización. La Ley del
Sector Eléctrico y la Ley de evaluación ambiental proporcionan el marco
general de estas materias
\cite{jefaturadelestado_ley242013_2013,jefaturadelestado_ley212013_2013}.
La descripción siguiente introduce los conceptos necesarios para interpretar
las publicaciones; no pretende recoger todas las variantes normativas ni
representar una secuencia obligatoria e idéntica para cada instalación.

La tabla~\ref{tab:siglas-administrativas} reúne las siglas que se utilizarán
en esta explicación. Conviene distinguir el procedimiento ambiental, el
estudio que aporta el promotor y la decisión del órgano ambiental.
\begin{table}[htbp]
\centering
\small
\renewcommand{\arraystretch}{1.2}
\begin{tabular}{L{1.4cm} L{11.5cm}}
\hline
Sigla & Significado en esta memoria \\
\hline
PACC & Permisos de acceso y conexión a la red eléctrica \\
EIA & Evaluación de impacto ambiental: procedimiento de evaluación \\
EsIA & Estudio de impacto ambiental: documento preparado por el promotor \\
DIA & Declaración de impacto ambiental: pronunciamiento del órgano ambiental \\
AAP & Autorización administrativa previa \\
AAC & Autorización administrativa de construcción \\
DUP & Declaración de utilidad pública \\
AE & Autorización de explotación \\
\hline
\end{tabular}
\caption[Siglas administrativas]{Siglas para distinguir permisos,
procedimientos, documentos y decisiones. Las referencias normativas y sus
límites se indican en el texto.}
\label{tab:siglas-administrativas}
\end{table}

\subsection{Acceso y conexión a la red}
Los permisos de acceso y conexión se refieren a la posibilidad de utilizar
la red eléctrica y a las condiciones para conectar la instalación. Son
conceptos distintos de la autorización para construirla o explotarla.
En este TFM aportan contexto para comprender la tramitación, pero no se
reconstruye un registro completo de permisos de red.
\citapendiente{Verificar la normativa específica de acceso y conexión y
los preceptos aplicables; la bibliografía actual solo proporciona el marco
general del sector eléctrico.}

\subsection{Evaluación ambiental e información pública}
La evaluación ambiental examina los efectos que puede tener un proyecto
sobre el medio ambiente. En el procedimiento ordinario, el promotor prepara
un \emph{Estudio de Impacto Ambiental}, que describe el proyecto, sus
alternativas, los efectos previstos y las medidas propuestas. El estudio no
es la decisión de la administración.

La \emph{información pública} permite consultar la documentación y presentar
alegaciones. Su anuncio puede referirse a la documentación ambiental y a
autorizaciones solicitadas. Que una autorización se someta a información
pública no significa que haya sido concedida.

La \emph{Declaración de Impacto Ambiental} (DIA) recoge la decisión del órgano
ambiental en la evaluación ordinaria, con el pronunciamiento y las condiciones
que correspondan. La evaluación simplificada utiliza un informe de impacto
ambiental. La diferencia entre estudio, trámite de consulta y decisión
ambiental es necesaria para no tratar sus menciones como hitos equivalentes
\cite{jefaturadelestado_ley212013_2013}.
\citapendiente{Completar las citas con los artículos y la versión normativa
consultada para estudio, información pública, DIA e informe de impacto
ambiental. No se han verificado textos legales consolidados en esta fase.}

\subsection{Autorizaciones y utilidad pública}
La \emph{autorización administrativa previa} (AAP) se refiere al proyecto
de la instalación en los términos previstos por la normativa eléctrica.
La \emph{autorización administrativa de construcción} (AAC) permite su
construcción conforme al proyecto técnico aprobado. Aunque puedan aparecer
en una misma resolución, representan actuaciones distintas.

La \emph{declaración de utilidad pública} (DUP) se relaciona con los bienes
y derechos afectados por la instalación, a efectos de expropiación o
servidumbre. Puede recaer sobre una infraestructura compartida, como una
línea de evacuación, lo que obliga a identificar el objeto de la decisión.
Finalmente, la \emph{autorización de explotación} se refiere a la puesta en
servicio de la instalación. La existencia de una AAP o una AAC no demuestra
por sí sola que la planta esté operativa.

El procedimiento de autorización de instalaciones eléctricas se desarrolla,
entre otras disposiciones, en el Real Decreto 1955/2000
\cite{ministeriodeeconomía_realdecreto1955_2001}.
\citapendiente{Verificar los preceptos de AAP, AAC, DUP y explotación,
sus modificaciones y los datos bibliográficos del Real Decreto 1955/2000;
la entrada existente utiliza el año 2001.}

Para el seguimiento documental interesa distinguir el \emph{tipo de
actuación} de su \emph{decisión}. Por ejemplo, una publicación puede someter
una AAP a información pública y otra conceder esa misma clase de
autorización. El sistema conserva ambos elementos, junto con la publicación
que respalda cada uno, sin completar los pasos ausentes por suposición.

La figura~\ref{fig:procedimiento-administrativo} representa un recorrido
orientativo con evaluación ambiental ordinaria y autorización. Las vías
ambiental y eléctrica pueden avanzar parcialmente en paralelo, aunque el
pronunciamiento de la DIA condiciona la autorización. La DUP se tramita cuando
procede y no ocupa necesariamente una fase posterior a la AAC.
\citapendiente{Verificar los preceptos que respaldan el paralelismo parcial,
el condicionamiento ambiental y la relación temporal de la DUP con AAP/AAC.}

La puesta en servicio puede requerir pruebas y autorizaciones de explotación
provisionales o definitivas, según el supuesto aplicable.
\citapendiente{Precisar los supuestos y la terminología normativa de pruebas,
autorización provisional y autorización definitiva de explotación.}
Esta distinción aporta contexto: el sistema de extracción utiliza el tipo
general \texttt{autorizacion\_explotacion}, sin dos categorías estructuradas
separadas para esas modalidades.

\begin{figure}[p]
\centering
% F01: esquema conceptual; referencias normativas pendientes en el capítulo 2.
\begin{tikzpicture}
\node[tfm accent, text width=4.4cm] (proyecto) at (0,0) {Proyecto renovable};
\node[tfm data, text width=5.8cm] (red) at (0,-1.2)
  {Permisos de acceso\\y conexión (PACC)};
\node[tfm heading, text width=4.5cm] at (-3.2,-2.15) {Procedimiento ambiental};
\node[tfm heading, text width=4.5cm] at (3.2,-2.15) {Procedimiento eléctrico};
\node[tfm box, text width=4.2cm] (estudio) at (-3.2,-3.35)
  {EIA / EsIA\\Evaluación y estudio};
\node[tfm box, text width=4.2cm] (consultas) at (-3.2,-4.65)
  {Información pública\\y consultas};
\node[tfm accent, text width=4.2cm] (dia) at (-3.2,-5.95) {DIA};
\node[tfm box, text width=4.2cm] (solicitudes) at (3.2,-3.35)
  {Solicitudes de autorización\\AAP / AAC};
\node[tfm box, text width=4.2cm] (publica) at (3.2,-4.65)
  {Tramitación e\\información pública};
\node[tfm accent, text width=4.2cm] (autorizaciones) at (3.2,-7.25)
  {Resoluciones de autorización\\AAP / AAC};
\node[tfm optional, text width=4.2cm] (dup) at (-3.2,-8.0)
  {DUP, cuando proceda\\Sin posición temporal fija\\respecto de la AAC};
\node[tfm box, text width=4.2cm] (obra) at (3.2,-8.65) {Construcción};
\node[tfm box, text width=4.2cm] (ae) at (3.2,-10.0)
  {Puesta en servicio: AE\\Pruebas, cuando correspondan};
\node[tfm accent, text width=4.2cm] (operacion) at (3.2,-11.35) {Operación};
\draw[tfm flow] (proyecto) -- (red);
\draw[tfm relation] (red.south) -- (0,-2.65);
\draw[tfm flow, dashed] (0,-2.65) -| (estudio.north);
\draw[tfm flow, dashed] (0,-2.65) -| (solicitudes.north);
\draw[tfm flow] (estudio) -- (consultas);
\draw[tfm flow] (consultas) -- (dia);
\draw[tfm flow] (solicitudes) -- (publica);
\draw[tfm flow] (publica) -- (autorizaciones);
\draw[tfm flow] (dia.east) -- (0,-5.95)
  node[tfm note, above, text width=2.0cm] {condiciona}
  |- (autorizaciones.west);
\draw[tfm relation] (dup.east) -- (0,-8.0) |- (autorizaciones.west);
\draw[tfm flow] (autorizaciones) -- (obra);
\draw[tfm flow] (obra) -- (ae);
\draw[tfm flow] (ae) -- (operacion);
\node[tfm note, text width=11.3cm] at (0,-12.6)
  {Esquema orientativo de un recorrido con autorización y evaluación ordinaria.\\
   Las vías pueden avanzar parcialmente en paralelo; no es un calendario universal.};
\end{tikzpicture}

\caption[Procedimiento administrativo simplificado]{Un proyecto puede
generar publicaciones sobre solicitudes, consultas y resoluciones diferentes.
La DIA condiciona la autorización; la DUP es condicional y su enlace discontinuo
no fija un orden respecto de la AAC. Se omiten denegaciones y otras variantes.
Elaboración propia; respaldo normativo pendiente en los TODO-CITA del texto.}
\label{fig:procedimiento-administrativo}
\end{figure}

\section{El BOE y su acceso programático}
\label{sec:boe-api-concepto}
El BOE es una fuente oficial de disposiciones, resoluciones y anuncios.
Ofrece sumarios diarios, que reúnen las publicaciones de cada fecha, y
documentos consultables individualmente. Localizar y descargar esos documentos
uno por uno exigiría repetir muchas operaciones manuales. Además, sería
necesario registrar qué fechas y publicaciones se han consultado para poder
repetir el proceso.

El \emph{acceso programático} permite que un programa solicite información
mediante instrucciones. Una interfaz de programación de aplicaciones, o
\emph{API}, define cómo realizar esas solicitudes y cómo interpretar las
respuestas. La web está pensada para que una persona consulte el BOE; la API
de sumarios permite que un programa solicite la relación de publicaciones de
una fecha de forma estructurada
\cite{agenciaestatalboletínoficialdelestado_apiparaacceso_2024}.

En este TFM la API se utiliza para obtener los sumarios y sus metadatos, es
decir, los datos que describen cada publicación: identificador, fecha, título
y enlaces, entre otros. Un filtro selecciona las publicaciones candidatas y,
después, el programa descarga sus documentos mediante los enlaces XML
proporcionados en el sumario. XML es un formato que organiza contenido
mediante etiquetas. El sumario ayuda a localizar el documento; el documento
aporta el texto que se analizará. Las solicitudes y el flujo concreto se
describen en la sección~\ref{sec:adquisicion-api}.

Cada publicación dispone de identificador BOE, fecha y enlaces de consulta.
Los sumarios aportan además título y contexto de clasificación, como sección
y departamento. El cuerpo documental puede incluir antecedentes, fundamentos,
una parte resolutiva y anexos, aunque los anuncios no siempre siguen esa
estructura. Estas diferencias afectan a la localización e interpretación de
la información.

La principal ventaja es la posibilidad de contrastar un dato con una
publicación identificable. Sin embargo, el identificador BOE identifica el
documento, no la planta. La información utilizada por este sistema no ofrece
una clave estable y común de proyecto entre publicaciones. Las denominaciones
pueden incluir abreviaciones o descriptores diferentes, y un anuncio puede
referirse a más de una planta.

Además, una historia restringida al BOE puede omitir actuaciones publicadas
en otros boletines o no recuperadas por la selección documental. La ausencia
de una publicación en el corpus no equivale a la ausencia del trámite.
\citapendiente{Añadir una referencia oficial sobre el alcance del BOE y la
distribución competencial y de publicaciones entre administraciones.}

\section{El problema del seguimiento longitudinal}
El seguimiento es \emph{longitudinal} cuando relaciona observaciones del
mismo proyecto en momentos diferentes. En este trabajo se parte de menciones
documentales, es decir, de las apariciones de una planta en cada publicación,
y se comprueba si pueden reunirse bajo una misma identidad.

Por ejemplo, Don Rodrigo II aparece en un anuncio de información pública y,
meses después, en una resolución que concede las autorizaciones solicitadas.
La tabla~\ref{tab:don-rodrigo} resume estas dos publicaciones del corpus final.
\footnote{Publicaciones de consulta:
\href{https://www.boe.es/diario_boe/txt.php?id=BOE-B-2026-12663}{BOE-B-2026-12663}
y \href{https://www.boe.es/diario_boe/txt.php?id=BOE-A-2026-17939}{BOE-A-2026-17939}.
El ejemplo procede de la auditoría temporal y de las correcciones aprobadas
del corpus W14.}
El anuncio de abril también incluye Bianor, que se conserva como proyecto
independiente. Compartir publicación no convierte las dos plantas en un
único proyecto.

\begin{table}[htbp]
\centering
\small
\renewcommand{\arraystretch}{1.2}
\begin{tabular}{L{2.0cm} L{4.0cm} L{6.3cm}}
\hline
Fecha & Publicación & Actuación actual sobre Don Rodrigo II \\
\hline
23/04/2026 & BOE-B-2026-12663 & Información pública de las solicitudes de AAP y AAC \\
19/08/2026 & BOE-A-2026-17939 & Concesión de AAP y AAC \\
\hline
\end{tabular}
\caption[Dos publicaciones de Don Rodrigo II]{Dos publicaciones de Don
Rodrigo II. La referencia a abril en la resolución de agosto es un antecedente,
no una nueva información pública. Fuente: auditoría temporal W14 y
correcciones aprobadas, con los enlaces BOE indicados en el texto.}
\label{tab:don-rodrigo}
\end{table}

La resolución de agosto recuerda la información pública de abril entre sus
antecedentes. Si esa referencia se registra como una nueva información
pública de agosto, la cronología duplica el trámite y lo sitúa en el momento
equivocado. La auditoría del corpus identificó esa confusión y la corrección
aprobada conservó en agosto las concesiones de AAP y AAC, excluyendo las
actuaciones que correspondían a antecedentes.

El ejemplo muestra dos decisiones diferentes: relacionar las menciones de
Don Rodrigo II entre documentos y atribuir cada actuación a la publicación
que realmente la comunica como actual. La metodología debe resolver ambas
sin confundir la fecha de publicación con todas las fechas citadas en el texto.
